\documentclass[11pt]{article}
\usepackage[a4paper,margin=26mm]{geometry}
\usepackage{amsmath,amssymb,amsthm,booktabs}
\usepackage[colorlinks=true,urlcolor=blue,linkcolor=blue]{hyperref}
\newtheorem{theorem}{Theorem}
\newtheorem{lemma}[theorem]{Lemma}
\newcommand{\OPT}{\operatorname{OPT}}
\newcommand{\cost}{\operatorname{cost}}
\newcommand{\MS}{\operatorname{MS}}
\newcommand{\Splay}{\operatorname{Splay}}
\newcommand{\lgit}[1]{\log^{(#1)}}
\setlength{\emergencystretch}{3em}
\title{Standard multi-splay after Cha-Cha:\\a common-forest comparison}
\author{Research session of 2 October 2026}
\date{}
\begin{document}
\maketitle

\begin{abstract}
The triple-iterated-logarithmic bound of the local Cha-Cha
manuscript does not transfer to standard multi-splay by replacing a
balanced auxiliary with a splay tree. We obtain a weaker transfer:
subject to the two local manuscripts' core proofs, standard multi-splay
has competitive factor
$O((1+\lgit{3}n)^2(1+\lgit{4}n)^2)$, with the usual additive
initialization allowance. The argument uses the exact mixed split/join
word executed by multi-splay, a persistent forest extension of the
current splay comparison, and a compatible Cha-Cha comparator.
All new composition and endpoint arguments here are hand proofs;
the existing Lean results certify specified constituent interfaces,
not this new conclusion. The exact $O(\lgit{3}n)$ bound for standard
multi-splay remains unproved in this session.
\end{abstract}

\paragraph{Notation and scope.}
$\lgit{j}$ denotes the $j$-fold iterated logarithm, not a power of a
logarithm. All small logarithms are stabilized at a positive constant.
The algorithm under discussion is the standard Wang--Derryberry--Sleator
multi-splay algorithm, including its final switch at the requested node.
Its fixed reference is the usual minimum-height node-storing BST.
Set $H=\lceil\log_2(n+1)\rceil$ and
\[
 \lambda=1+\log_2\log_2(H+2),\qquad
 g(H)=(1+\log\log(H+4))(1+\log\log\log(H+16))^2.
\]
Thus $\lambda=\Theta(1+\lgit{3}n)$ and
$g(H)=O((1+\lgit{3}n)(1+\lgit{4}n)^2)$.
Costs are pointer moves and rotations in the BST model. The comparator
may compute from already read information for free, as in the local
Cha-Cha manuscript. No uniform word-RAM implementation is asserted.

\section{What the directory supplies, and what it does not}

The splay manuscript supplies
\[
 \cost_{\Splay}(X)=
 O\bigl((\OPT(X)+n)(1+\lgit{2}n)(1+\lgit{3}n)^2\bigr).
                                                        \tag{1}
\]
Its split appendix and the existing heavy-path and rib representation
interfaces are especially relevant here. The newer historical-compression
notes improve the smaller factors but do not by themselves remove the
leading $\lgit{2}n$ factor. We use the simpler accounting underlying (1).

The current unified Cha-Cha manuscript in the neighboring
repository is newer than the splay-side bridge notes. It presents a
complete hand upper proof for its specific nested representation.
Its lower theorem is formally checked in the actual BST model. With
primary alternations $P$, completed episodes $D$, and initialized
fixed-owner secondary changes $F$, it gives
\[
 m+P+D\le2888\OPT(X),\qquad
 F\le2D+2P+5(n-1),
\]
and its algorithm costs $O((m+P+F+n)\lambda)$.

The physical upper proof retains a secondary decomposition of each
primary path. Its smallest balanced auxiliaries have only
$O(\log H)$ nodes. Standard multi-splay instead stores all $O(H)$
nodes of a primary path in one splay auxiliary. Applying its ordinary
access lemma directly still costs $O(\log H)=O(\lgit{2}n)$ per switch.
The stronger lower counter alone does not reduce that upper bound.

\section{An endpoint-aware lower counter for standard multi-splay}

For request $x_i$, let
\[
 E_i=(d_1<\cdots<d_k=\operatorname{depth}_P(x_i))
\]
consist of the wrong proper-ancestor depths and the mandatory final
switch depth. Write $S=\sum_i|E_i|$. At a reference leaf the final
switch still performs its splays; an analytical null-child direction
bit may be toggled without changing the physical algorithm.

On a balanced binary interval hierarchy $B$ of depths $0,\ldots,H$,
an emission $d$ writes left or right at every interval containing $d$.
The row is indexed by the interval and the original request's ancestor
at that interval's left endpoint. All row bits start right. Let
$F_\uparrow$ use the increasing order $E_i$, and let $F_{\rm act}$
use the actual order
\[
 (d_{k-1},\ldots,d_1,d_k).
\]
There is no extra completion emission $H$.

\begin{lemma}[Endpoint and order transport]\label{lem:counter}
For the standard reference,
\[
 m+S+F_\uparrow=O(\OPT(X)+n),\qquad
 |F_{\rm act}-F_\uparrow|
 \le2\sum_i\max(|E_i|-2,0).
\]
Consequently $m+S+F_{\rm act}=O(\OPT(X)+n)$.
\end{lemma}

\begin{proof}[Proof outline]
Replace each original key $q$ at depth $d$ by a pivot at depth $2d$
and two connector children. The left connector separates the original
left subtree from a copy $q^-$; the right connector separates a copy
$q^+$ from the original right subtree. Query the copy on the side
opposite the current preferred child of $q$. The even-depth pivot
changes are then exactly the ancestor switches and the final toggle.
Pad copy leaves by all-right paths to a common depth. Only the $3n$
pivots and connectors can change; the padding creates no primary flips.

Replace each secondary depth leaf $d$ by the pair $2d,2d+1$.
The original row word is a subsequence of the corresponding expanded
row word. Initialized variation cannot increase on deletion.
The lower-bound proof uses ordered binary interval partitions and
unique separation of two depths; balancing of the secondary hierarchy
is not needed for that lower proof.

The apparently large padded universe causes no large initial credit.
At expanded depth $a$ there are at most
$2^{\lfloor a/2\rfloor+1}$ core owners. Distinct secondary splits use
distinct positions, so the number of rows ever written left is at most
$4(2^H-1)=O(n)$. A competing BST on the original keys serves either
copy using a two-key gadget per original node, at constant overhead.
The Cha-Cha lower theorem therefore proves the first assertion.

For the order assertion, fix one query and one secondary interval.
A changed mixed block has its initialized variation altered by at
most two, including its entry and exit boundaries. The induced binary
trie on $k$ emitted depths has $k-1$ branching nodes. The parent of
the maximum emitted leaf in the compressed trie has only that maximum
on its right side, and its block is unchanged
when the maximum stays last. At most $k-2$ mixed blocks change.
Sum this replacement bound over queries and fixed owner rows.
\end{proof}

The complete proof and the initialization constants are in
\href{endpoint-embedding.pdf}{\texttt{endpoint-embedding.pdf}}.
An exhaustive finite audit in \texttt{endpoint-check.py} checked the
encoding, row words, order estimate, and initialization support on
32,140 sequences containing 171,611 requests. This is a check of those
identities, not a substitute for the lower theorem or an upper-cost test.

\section{The exact mixed forest word}

A forest component is one auxiliary BST, with unit-cost selection of
its root shared by both executions. No interior pointer jump is free.
The operations are root-splaying searches; a split which splays and
removes its pivot; and a raw common root join $x(L,R)$.

Consider a left-to-right switch at $y$. Its old component is
$U\cup\{y\}\cup L$, where $U$ is the preferred prefix above $y$ and
$L$ is the old left tail. The incoming right tail $R$ is another
component. Let $\ell,y,r$ be the nearest lower ancestor, the switched
key, and the nearest upper ancestor. The exact word, when both guards
exist, is
\[
 \begin{split}
 &\operatorname{Split}(y),\quad
  \operatorname{Split}(\ell),\quad
  \operatorname{Split}(r),\\
 &\operatorname{Join}(A,\ell,\varnothing),\quad
  \operatorname{Join}(R,r,B),\quad
  \operatorname{Join}(A',y,B').
 \end{split}                                             \tag{2}
\]
The guard splits occur in the two daughters of the first split.
Their detached pivots are then restored by the displayed joins.
The old tail $L$ remains parked. Missing guards omit their operations;
right-to-left switches are reflected.

The three split splays are exactly the full and stopped splays of
standard multi-splay, and the joins restore exactly its auxiliary
shapes. Foreign marked subtrees remain opaque in the gaps. Thus one
shape-independent word $W_{\MS}(X)$ describes the whole algorithm,
with at most six operations per switch and component size at most $H$.
The standard multi-splay traversal argument pays its initial global
search from these splays and the final switch.

\section{Splay comparison over one persistent forest}

\begin{theorem}[Common-forest extension]\label{thm:forest}
Assuming the current splay manuscript's core accounting proof, a
mixed forest word with $N$ initially active keys, $q$ structural
operations, maximum component size $K$, and a comparator of cost $C$
has splay-forest cost
\[
 O\bigl((C+N+q)g(K)\bigr).
\]
There is one global initialization allowance, not one per component life.
\end{theorem}

\begin{proof}[Proof outline and interface audit]
The existing Lean heavy representation has root split and join
interfaces with a uniform supported-path cap $O(\log^2 K)$ and
constant amortized allowance. Its potential is at most $64$ times
the number of real keys. The rib representation has the corresponding
interfaces, constant-cost amortized reorientation, and potential at
most $256$ times the number of real keys. Sum each potential over
live components. Root splits return the two actual child histories;
root joins consume those histories. This yields a comparison forest
of cost $O(C+N+q)$, height $O(\log^3 K)$, bounded minor ribs, and
rotations at constant depth.

At a common raw join, give the new pivot comparison rank zero and
increase every old child rank by one. All internal rank differences,
heavy paths, and lazy-interval histories of each child stay unchanged.
Only the two former child roots become new heap children. Each has
gap one. In the actual interface their contracted value changes from
one to two, requiring one unit value increase and one legal insertion
into a new singleton growing interval. Thus a join creates only
constantly many paid events and bends. Root deletion uses the split
appendix's inverse operation. Inactive keys can retain baseline value
one, or one may bound the total key universe by $N+q$.

The original pairing budgets now sum over the forest. Writing
$B_0=C+N+q$ and $G=O(1+\log\log K)$, the interval, boundary,
value-change, and bend budgets have exactly the original form with
$B_0$ as their input allowance. The same absorption step gives
$O(B_0G(1+\log G)^2)=O(B_0g(K))$ total work.
\end{proof}

This is a hand extension, with a concrete trace through the checked
local interfaces. It is not asserted to follow from the fixed-universe
statement (1) alone. The full interface account, source module names,
and exact word (2) are in
\href{common-forest-normal-form.pdf}{\texttt{common-forest-normal-form.pdf}}.

\section{A compatible Cha-Cha comparator}

\begin{theorem}[The comparator for the same word]\label{thm:comparator}
The exact word $W_{\MS}(X)$ admits a comparator of cost
\[
 C_{\rm forest}=O\bigl(n+\lambda(S+F_{\rm act})\bigr).
                                                               \tag{3}
\]
\end{theorem}

\begin{proof}[Proof outline]
View the node reference as a routing tree on its external gaps. Each
primary path is represented by the Cha-Cha depth-ordered tree,
unfolded onto exactly its own original separator keys. For a path
beginning at depth $a$, use genuine secondary intervals starting at
depth at least $a$, together with the structural boundary spine
$\beta_a$ of the induced suffix $[a,H]$.

The forest word already supplies the component and switched pivot.
Therefore the comparator can follow the actual bottom-up order without
paying for a separate discovery search. A local request to depth $d$
costs $O((1+f)\lambda)$, where $f$ counts genuine changes. Remove
the one administrative seam at the end of $\beta_a$, perform the inner
Tango request and the exposed primary splice, then restore the at most
one temporary right boundary preference and the seam. These are
constantly many extra balanced split/concatenate operations.

For every fixed genuine owner row, the comparator's write word is the
full $F_{\rm act}$ word with writes deleted whenever the selected
primary component begins below that owner. Splices preserve the
surviving histories. Hence the total genuine changes are at most
$F_{\rm act}$ by initialized variation monotonicity.

The remaining requirement is to obey the exact pivot checkpoints in
(2). After the local depth request and the conceptual tail cut, the
last two prefix gaps are cheaply exposed. In physical key order their
boundaries are $\ell,y$ and $y,r$. Both neighboring keys of an external
BST gap lie on its search path. A rooted hull of $O(\lambda)$ nodes
therefore contains all existing guards and $y$. During this preparation
the old tail remains attached inside the original common component.

Keep the rooted touched hulls of the ideal exposed-splice transaction
and its incoming component. Their untouched frontier subtrees are
identical ordered key intervals before and after. Reconfigure the
old hull to $y(\ell(A,L),r(\varnothing,B))$, execute the prescribed
splits and joins, and obtain $y(\ell(A,\varnothing),r(R,B))$ with
$L$ parked. Reconfigure each resulting rooted hull separately to its
valid Cha-Cha after-state. Every rotation occurs in a current
common component. The total touched-hull cost is
$O((1+f)\lambda)$; no frontier subtree is rebuilt.

Summing proves (3). Initially every physical component can be put
into its prescribed balanced representation in linear cost, totaling
$O(n)$. For a nonperfect minimum-height reference, complete only the
virtual routing reference and retain the all-right descendant of
each original external gap. Inactive pieces and dummy separators have
no physical images. This preserves the exact original word, all real
guards, and the cost bound without introducing physical dummy keys.
\end{proof}

The full proof, including the suffix seam and the checkpoint certificate,
is in \href{chacha-common-forest-comparator.pdf}{\texttt{chacha-common-forest-comparator.pdf}}.

\section{The resulting bound and the precise remaining question}

\begin{theorem}[Composed hand upper bound]\label{thm:main}
Under the stated local Splay and Cha-Cha proof dependencies,
standard multi-splay on the usual minimum-height reference satisfies
\[
 \boxed{\displaystyle
 \cost_{\MS}(X)=
 O\bigl((\OPT(X)+n)(1+\lgit{3}n)^2(1+\lgit{4}n)^2\bigr).}
\]
\end{theorem}

\begin{proof}
Lemma~\ref{lem:counter} and Theorem~\ref{thm:comparator} give
$C_{\rm forest}=O(\lambda(\OPT+n))$.
Apply Theorem~\ref{thm:forest} to the exact word, with $K=H$,
$N=n$, and $q\le6S$. The result is
$O((\OPT+n)\lambda g(H))$, which is the displayed expression.
\end{proof}

There are two different losses in this derivation: the compatible
balanced comparator pays $\lambda$, and the current splay comparison
pays $g(H)$. They multiply. Replacing this product by $\lambda$ needs
an additional upper-bound argument about the actual multi-splay
rotations. One sufficient target is a single persistent potential with
\[
 \cost_{\MS}(i)+\Phi_i-\Phi_{i-1}
 \le C\lambda\bigl(1+s_i+f_i\bigr),\qquad
 \Phi_0-\inf_i\Phi_i=O(n\lambda),                       \tag{4}
\]
where $s_i,f_i$ are the exact standard-switch and owner-row changes.
Equation (4) is unproved here.

A weaker useful intermediate target is an $O(\OPT+n)$ comparator
for the same entire forest word. Combined with Theorem~\ref{thm:forest}
it would remove one factor $\lambda$, leaving
$O((1+\lgit{3}n)(1+\lgit{4}n)^2)$. Independently optimal trees for
individual path lifetimes do not supply such a comparator: their initial
shapes and their boundary exposures must be transported through (2).

\paragraph{Why several shortcuts fail.}
Restarting (1) whenever a preferred path is born can charge
$\Omega(mH)$ for repeated requests at the reference root, although
their actual cost is $O(m)$ after initialization. A sum of functions of
component sizes cannot repair this in a perfect reference: its path-size
multiset is invariant. The classical forest sum-log-size potential has
$\Theta(\log H)$ positive jumps even on constant-work exposed switches.
These are obstructions to specific proof strategies, not lower bounds
on multi-splay's competitive ratio. Details are in
\href{order-and-potential.pdf}{\texttt{order-and-potential.pdf}} and
\href{ClassicalPotentialObstruction.pdf}{\texttt{ClassicalPotentialObstruction.pdf}}.

\paragraph{A modified algorithm does achieve the desired factor.}
If one retains Cha-Cha's secondary preferred-path decomposition
and replaces only its $O(\log H)$-sized balanced auxiliaries by splay
trees, one global virtual sum-log-size potential proves
$O(\lambda(\OPT+n))$. The cut keeps each child's state, the role
exchange leaves the virtual potential unchanged, and the final graft
increases it by at most $O(\log\log H)$. This is proved in
\href{NestedSplayAuxiliaries.pdf}{\texttt{NestedSplayAuxiliaries.pdf}}.
It is a different algorithm from standard multi-splay: the extra
secondary decomposition is part of the construction.

\section{Verification status and sources}

The new endpoint embedding, order comparison, mixed forest account,
compatible comparator, and modified nested algorithm were independently
hand-audited during this session. The finite checks cover only explicitly
listed combinatorial identities and potential calculations. No new
Lean theorem asserting the composed multi-splay bound has been added.
The precise existing checked dependencies are named in the detailed
forest note; the arbitrary-secondary-hierarchy lower extension and
the active-gap padding are hand arguments.

The underlying manuscripts and primary literature are:
\begin{itemize}
\item The current local splay manuscript, \path{../splay/main.tex}, its
split appendix, and \path{../splay/formalization/Splay/Accounting/}.
The public main bound is in
\href{https://arxiv.org/abs/2607.18498}{\emph{Splay trees are almost dynamically optimal}}.
\item The current neighboring manuscript, \texttt{chacha.tex},
in \path{../chacha/research/}.
\item C.~C.~Wang, J.~Derryberry, and D.~D.~Sleator,
\href{https://www.cs.cmu.edu/~jonderry/multi-splay.pdf}{\emph{$O(\log\log n)$-Competitive Dynamic Binary Search Trees}}, SODA 2006;
and \href{https://www.csd.cmu.edu/sites/default/files/phd-thesis/CMU-CS-06-140.pdf}{Wang's thesis}, especially the exact switch and traversal conventions.
\item S.-L.~Woo,
\href{https://csd.cs.cmu.edu/sites/default/files/phd-thesis/CMU-CS-09-133.pdf}{\emph{Heterogeneous Decomposition of Degree-Balanced Search Trees and Its Applications}}, 2009.
Its conditional multi-splay dynamic-finger corollary expressly includes
splits and joins in the required splay comparison.
\end{itemize}

\end{document}
